\begin{table*}[tb]
    \centering
    \resizebox{0.99\textwidth}{!}{%
    \begin{tabular}{|c|c|ccccc|c|c}
        \Xhline{1.5pt}
        \textbf{Input Type} & \textbf{Model} & \textbf{One-hop} & \textbf{Conjunction} & \textbf{Existence} & \textbf{Multi-hop} & \textbf{Negation} & \textbf{Total} \\ \hline
        \multirow{2}{*}{\textit{Claim Only}}  
        & \textsc{FactKG} BERT Baseline & 69.64 & 63.31 & 61.84 & 70.06 & 63.62 & 65.20 \\
        & FactGenius RoBERTa Baseline & 71 & 72 & 52 & 74 & 54 & 68 \\
        & BERT (no subgraphs) & 67.71 & 67.48 & 62.51 & 73.28 & 64.23 & 68.99 \\
         \hline
        \multirow{3}{*}{\textit{With Subgraphs}}
        & \textsc{FactKG} GEAR Benchmark & 83.23 & 77.68 & 81.61 & 68.84 & 79.41 & 77.65 \\
        & FactGenius RoBERTa-two-stage & 89 & 85 & 95 & 75 & 87 & 85 \\
        & QA-GNN (single-step) & 79.08 & 74.43 & 83.37 & 74.72 & 79.60 & 78.08 \\
        & BERT (single-step) & \textbf{97.40} & \textbf{97.51} & \textbf{97.31} & \textbf{80.32} & \textbf{92.54} & \textbf{93.49} \\
         \hline
    
        \Xhline{1.5pt}
    \end{tabular}
    }
    \caption{\textbf{Test-set accuracy for the best models from this article and the best benchmark models.} The \textsc{FactKG} models are from \citet{kim2023factkg}, while the FactGenius models are from \citet{gautam2024factgenius}. The fine-tuned BERT model performed the best, while the QA-GNN was the computationally most efficient model.}
    \label{tab:best_test_accuracies}
\end{table*}